% --- scholar LaTeX fragment --------------------------------------------
% Generated by: scholar sessions export -f table
% Session: t-interaction-brittle
% \input this file from the body of a document whose
% preamble loads:
%   \usepackage{longtable}
%   \usepackage{booktabs}
%   \usepackage{array}
% ----------------------------------------------------------------------

\begingroup\footnotesize
\begin{longtable}{@{}%
  >{\raggedright\arraybackslash}p{0.44\textwidth}c%
  >{\raggedright\arraybackslash}p{0.13\textwidth}%
  >{\raggedright\arraybackslash}p{0.26\textwidth}@{}}
\caption{Records bearing on the claim (5 cited, 0 supporting and 0 qualifying or contradicting it, out of 123 unique hits; the 90 adjacent and 28 false hits are not listed). The last column gives the reason for each inclusion; for machine-classified rows the model's confidence is given. Databases: OpenAlex (OA; not open access; incomplete for 3 of 3 queries), DBLP (incomplete for 3 of 3 queries), IEEE Xplore (IEEE; incomplete for 3 of 3 queries), Web of Science (WoS), Scopus, arXiv (incomplete for 3 of 3 queries).}\label{tab:t-interaction-brittle}\\
\toprule
Title & Year & Provider & Reason \\
\midrule
\endfirsthead
\multicolumn{4}{@{}l}{\emph{\tablename~\thetable{} (cont.)}}\\
\toprule
Title & Year & Provider & Reason \\
\midrule
\endhead
\midrule \multicolumn{4}{r@{}}{\emph{continued on next page}}\\
\endfoot
\bottomrule
\endlastfoot
Mock objects for testing java systems & 2018 & DBLP, IEEE, OA, arXiv & \emph{cited}: mock-coupling \\
Point/counterpoint & 2007 & Scopus & \emph{cited}: practitioner-opinion \\
To Mock or Not to Mock? An Empirical Study on Mocking Practices & 2017 & DBLP, IEEE, OA, arXiv & \emph{cited}: mock-coupling \\
Toward a Resilience-Oriented Understanding of Unit Test Suites and Refactoring in Software Evolution & 2025 & Scopus & \emph{cited}: fragile-tests \\
Understanding and Characterizing Mock Assertions in Unit Tests & 2025 & DBLP, IEEE, OA, arXiv & \emph{cited}: mock-assertion-value \\
\end{longtable}
\endgroup
